% ============================================================
% 数学\数据表格\三年级.tex
% 本文件由 main.tex 通过 \input 引入，请勿单独编译
% ============================================================


\section{解答：先把下面的数位顺序表填写完整，再填一填}

\vspace{0.4cm}

\begin{center}
\renewcommand{\arraystretch}{1.5}
\begin{tabular}{|c|c|c|c|c|c|c|c|c|c|c|c|c|}
    \hline
    & \multicolumn{6}{c|}{整数部分} 
    & 小数点 
    & \multicolumn{5}{c|}{小数部分} \\
    \hline
    
    数位 
    & $\cdots$ 
    & \textbf{\textcolor{red!70!black}{\makecell{万\\位}}} 
    & \textbf{\textcolor{red!70!black}{\makecell{千\\位}}} 
    & \textbf{\textcolor{red!70!black}{\makecell{百\\位}}} 
    & \textbf{\textcolor{red!70!black}{\makecell{十\\位}}} 
    & \makecell{个\\位} 
    & \multirow{2}{*}{$\cdot$}
    & \makecell{十\\分\\位} 
    & %\textbf{\textcolor{red!70!black}{\makecell{百\\分\\位}}}
    & \textbf{\textcolor{red!70!black}{\makecell{千\\分\\位}}} 
    & \textbf{\textcolor{red!70!black}{\makecell{万\\分\\位}}} 
    & $\cdots$ \\
    \cline{1-7} \cline{9-13}
    
    \makecell{计\\数\\单\\位} 
    & $\cdots$ 
    & \textbf{\textcolor{red!70!black}{万}} 
    & \textbf{\textcolor{red!70!black}{千}} 
    & \textbf{\textcolor{red!70!black}{百}} 
    & \textbf{\textcolor{red!70!black}{十}} 
    & \makecell{一\\(个)} 
    &
    & \textbf{\textcolor{red!70!black}{\makecell{十\\分\\之\\一}}} 
    & \textbf{\textcolor{red!70!black}{\makecell{百\\分\\之\\一}}} 
    & \textbf{\textcolor{red!70!black}{\makecell{千\\分\\之\\一}}} 
    & %\textbf{\textcolor{red!70!black}{\makecell{万\\分\\之\\一}}}
    & $\cdots$ \\
    \hline
\end{tabular}
\end{center}

\vspace{0.5cm}

\noindent
每相邻两个计数单位间的进率都是（\textbf{\textcolor{red!70!black}{10}}）。



%% ==============================
%% 操作题：无人机最短路线（从起飞点到对岸作垂线）
%% ==============================
\newpage




% ============================================================
% 来源：课课练·三下数学 作图题答案.tex
% ============================================================


\newpage

\begin{center}
  {\large\bfseries 解答：经过每两个点最多可以画几条直线}
\end{center}

\vspace{0.4cm}

\noindent
\begin{tikzpicture}[scale=1.0, line width=0.55pt]
  \pgfmathsetmacro{\r}{0.85}   % 各小图点圆半径

  %% ---- n=2：1条 ----
  \begin{scope}[xshift=0cm, yshift=0cm]
    \foreach \i in {1,2} {
      \pgfmathsetmacro{\ang}{90 + 180*(\i-1)}
      \coordinate (A\i) at (\ang:\r);
    }
    \draw (A1)--(A2);
    \foreach \i in {1,2} { \node[vertex dot, minimum size=4pt] at (A\i) {}; }
    \node[above, font=\footnotesize, gray]         at (0,  \r+0.25) {$2$ 个点};
    \node[below, font=\small\bfseries, red!70!black] at (0, -\r-0.35) {$1$ 条};
  \end{scope}

  %% ---- n=3：3条 ----
  \begin{scope}[xshift=3cm, yshift=0cm]
    \foreach \i in {1,2,3} {
      \pgfmathsetmacro{\ang}{90 + 120*(\i-1)}
      \coordinate (B\i) at (\ang:\r);
    }
    \foreach \i in {1,2,3} {
      \foreach \j in {1,2,3} {
        \pgfmathparse{int(\i < \j)}
        \ifnum\pgfmathresult=1 \draw (B\i)--(B\j); \fi
      }
    }
    \foreach \i in {1,2,3} { \node[vertex dot, minimum size=4pt] at (B\i) {}; }
    \node[above, font=\footnotesize, gray]           at (0,  \r+0.25) {$3$ 个点};
    \node[below, font=\small\bfseries, red!70!black] at (0, -\r-0.35) {$3$ 条};
  \end{scope}

  %% ---- n=4：6条 ----
  \begin{scope}[xshift=6cm, yshift=0cm]
    \foreach \i in {1,2,3,4} {
      \pgfmathsetmacro{\ang}{90 + 90*(\i-1)}
      \coordinate (C\i) at (\ang:\r);
    }
    \foreach \i in {1,2,3,4} {
      \foreach \j in {1,2,3,4} {
        \pgfmathparse{int(\i < \j)}
        \ifnum\pgfmathresult=1 \draw (C\i)--(C\j); \fi
      }
    }
    \foreach \i in {1,2,3,4} { \node[vertex dot, minimum size=4pt] at (C\i) {}; }
    \node[above, font=\footnotesize, gray]           at (0,  \r+0.25) {$4$ 个点};
    \node[below, font=\small\bfseries, red!70!black] at (0, -\r-0.35) {$\mathbf{6}$ 条};
  \end{scope}

  %% ---- n=5：10条 ----
  \begin{scope}[xshift=0cm, yshift=-3.5cm]
    \foreach \i in {1,2,3,4,5} {
      \pgfmathsetmacro{\ang}{90 + 72*(\i-1)}
      \coordinate (D\i) at (\ang:\r);
    }
    \foreach \i in {1,2,3,4,5} {
      \foreach \j in {1,2,3,4,5} {
        \pgfmathparse{int(\i < \j)}
        \ifnum\pgfmathresult=1 \draw (D\i)--(D\j); \fi
      }
    }
    \foreach \i in {1,2,3,4,5} { \node[vertex dot, minimum size=4pt] at (D\i) {}; }
    \node[above, font=\footnotesize, gray]           at (0,  \r+0.25) {$5$ 个点};
    \node[below, font=\small\bfseries, red!70!black] at (0, -\r-0.35) {$\mathbf{10}$ 条};
  \end{scope}

  %% ---- n=6：15条 ----
  \begin{scope}[xshift=3cm, yshift=-3.5cm]
    \foreach \i in {1,2,3,4,5,6} {
      \pgfmathsetmacro{\ang}{90 + 60*(\i-1)}
      \coordinate (E\i) at (\ang:\r);
    }
    \foreach \i in {1,2,3,4,5,6} {
      \foreach \j in {1,2,3,4,5,6} {
        \pgfmathparse{int(\i < \j)}
        \ifnum\pgfmathresult=1 \draw (E\i)--(E\j); \fi
      }
    }
    \foreach \i in {1,2,3,4,5,6} { \node[vertex dot, minimum size=4pt] at (E\i) {}; }
    \node[above, font=\footnotesize, gray]           at (0,  \r+0.25) {$6$ 个点};
    \node[below, font=\small\bfseries, red!70!black] at (0, -\r-0.35) {$\mathbf{15}$ 条};
  \end{scope}

  %% ---- n=7：21条 ----
  \begin{scope}[xshift=6cm, yshift=-3.5cm]
    \foreach \i in {1,2,3,4,5,6,7} {
      \pgfmathsetmacro{\ang}{90 + 360/7*(\i-1)}
      \coordinate (F\i) at (\ang:\r);
    }
    \foreach \i in {1,2,3,4,5,6,7} {
      \foreach \j in {1,2,3,4,5,6,7} {
        \pgfmathparse{int(\i < \j)}
        \ifnum\pgfmathresult=1 \draw (F\i)--(F\j); \fi
      }
    }
    \foreach \i in {1,2,3,4,5,6,7} { \node[vertex dot, minimum size=4pt] at (F\i) {}; }
    \node[above, font=\footnotesize, gray]           at (0,  \r+0.25) {$7$ 个点};
    \node[below, font=\small\bfseries, red!70!black] at (0, -\r-0.35) {$\mathbf{21}$ 条};
  \end{scope}

\end{tikzpicture}

\vspace{0.8cm}

% 完整答案表格
%\begin{center}
\renewcommand{\arraystretch}{1.6}
\begin{tabular}{|c|c|c|c|c|c|c|}
\hline
点数/个   & $2$ & $3$ & $4$ & $5$ & $6$ & $7$ \\
\hline
直线数/条 & $1$ & $3$
  & \textbf{\textcolor{red!70!black}{6}}
  & \textbf{\textcolor{red!70!black}{10}}
  & \textbf{\textcolor{red!70!black}{15}}
  & \textbf{\textcolor{red!70!black}{21}} \\
\hline
\end{tabular}
%\end{center}

\vspace{0.6cm}
\hrule
\vspace{0.3cm}

{\small\color{gray}
\textbf{规律：}$n$ 个点中每两点连一条直线，共可画
$\dfrac{n \times (n-1)}{2}$ 条。\\
例：$4$ 个点 $\to$ $\dfrac{4\times3}{2} = 6$ 条；
    $7$ 个点 $\to$ $\dfrac{7\times6}{2} = 21$ 条。
}


%% ==============================
%% 估一估，连一连
%% ==============================



% ============================================================
% 来源：竖排表格.tex
% ============================================================


\newpage

\begin{table}[htbp]
    \centering
    \renewcommand{\arraystretch}{1.5} % 增加行高，让表格显得不那么拥挤
    \begin{tabular}{|c|c|c|c|c|c|c|c|c|c|c|c|c|}
        \hline
        % 第一行
        & \multicolumn{6}{c|}{整数部分} 
        & 小数点 
        & \multicolumn{5}{c|}{小数部分} \\
        \hline
        
        % 第二行
        数位 
        & $\cdots$ 
        & \makecell{万\\位} 
        & \makecell{千\\位} 
        & \makecell{百\\位} 
        & \makecell{十\\位} 
        & \makecell{个\\位} 
        & \multirow{2}{*}{$\cdot$} % 合并两行作为小数点
        & \makecell{十\\分\\位} 
        & % 故意留空的百分位
        & \makecell{千\\分\\位} 
        & \makecell{万\\分\\位} 
        & $\cdots$ \\
        \cline{1-7} \cline{9-13} % 画横线，跳过第8列（小数点列）
        
        % 第三行
        \makecell{计\\数\\单\\位} 
        & $\cdots$ 
        & 万 & 千 & 百 & 十 
        & \makecell{一\\(个)} 
        & % 此处已被 multirow 占用
        & \makecell{十\\分\\之\\一} 
        & \makecell{百\\分\\之\\一} 
        & \makecell{千\\分\\之\\一} 
        & % 故意留空的万分之一
        & $\cdots$ \\
        \hline
    \end{tabular}
\end{table}

